\clearpage
\begin{landscape}
\begingroup
\fontsize{9.5}{12}\selectfont
\fancyhead[L]{\small NACHLASS. STÖRUNGEN DER PALLAS}
\renewcommand{\headrulewidth}{0pt}
\setlength{\lineskip}{4pt}\setlength{\lineskiplimit}{4pt}
\setlength{\arraycolsep}{4pt}
\pdfbookmark[2]{Allgemeine Störungen, Erste Rechnung, [4]}{pallas-V-fragment4}
\begin{center}[4.]\end{center}
\par\noindent\hspace*{\parindent}[Die Rechnungsvorschriften des vorigen Artikels sollen durch ein Beispiel erläutert werden, wobei \(\mu=48\), und die Werthe der \(A_0, A_1, A_2\ldots A_{\mu-1}\) vorgegeben seien.]
\begin{center}
\textit{Bestimmung der Coefficienten einer periodischen Function aus 48 Werthen derselben.}\par
\medskip
Die Function ist das \(T\) bei den Störungen der Pallas durch Jupiter\par
Mittlere Anomalie der Pallas \(=270^0\).
\end{center}
\[
\begin{array}{@{}r|r|r|r|r|r|r@{}}
A_0=+8{,}589&A_1=+10{,}595&A_2=+12{,}182&A_3=+13{,}009&A_4=+12{,}524&A_5=+9{,}838&A_6=+3{,}448\\
&A_{47}=+6{,}401&A_{46}=+4{,}110&A_{45}=+1{,}877&A_{44}=-0{,}255&A_{43}=-2{,}241&A_{42}=-4{,}051
\end{array}
\]
\[
\begin{array}{@{}r|r|r|r|r|r@{}}
A_7=-8{,}908&A_8=-30{,}028&A_9=-60{,}693&A_{10}=-92{,}504&A_{11}=-105{,}044&A_{12}=-86{,}954\\
A_{41}=-5{,}665&A_{40}=-7{,}066&A_{39}=-8{,}245&A_{38}=-9{,}195&A_{37}=-9{,}910&A_{36}=-10{,}388
\end{array}
\]
\[
\begin{array}{@{}r|r|r|r|r|r@{}}
A_{13}=-53{,}580&A_{14}=-24{,}130&A_{15}=-4{,}733&A_{16}=+6{,}187&A_{17}=+11{,}526&A_{18}=+13{,}464\\
A_{35}=-10{,}623&A_{34}=-10{,}612&A_{33}=-10{,}342&A_{32}=-9{,}838&A_{31}=-9{,}070&A_{30}=-8{,}047
\end{array}
\]
\[
\begin{array}{@{}r|r|r|r|r|r@{}}
A_{19}=+13{,}381&A_{20}=+12{,}114&A_{21}=+10{,}177&A_{22}=+7{,}891&A_{23}=+5{,}465&A_{24}=+3{,}036\\
A_{29}=-6{,}770&A_{28}=-5{,}243&A_{27}=-3{,}476&A_{26}=-1{,}487&A_{25}=+0{,}698&
\end{array}
\]
\par\noindent\hspace*{\parindent}Die Cosinus (2 fach):
\[
\begin{array}{@{}r|r|r|r|r|r|r@{}}
B_0=+17{,}178&B_1=+16{,}996&B_2=+16{,}292&B_3=+14{,}886&B_4=+12{,}269&B_5=+7{,}597&B_6=-0{,}603\\
B_{24}=+6{,}072&B_{23}=+6{,}163&B_{22}=+6{,}404&B_{21}=+6{,}701&B_{20}=+6{,}871&B_{19}=+6{,}611&B_{18}=+5{,}417
\end{array}
\]
\[
\begin{array}{@{}r|r|r|r|r|r@{}}
B_7=-14{,}573&B_8=-37{,}094&B_9=-68{,}938&B_{10}=-101{,}699&B_{11}=-114{,}954&B_{12}=-97{,}342\\
B_{17}=+2{,}456&B_{16}=-3{,}651&B_{15}=-15{,}075&B_{14}=-34{,}742&B_{13}=-64{,}203&
\end{array}
\]
\par\noindent\hspace*{\parindent}Die Cosinus 2 (4 fach):
\[
\begin{array}{@{}r|r|r|r|r|r|r@{}}
F_0=+23{,}250&F_1=+23{,}159&F_2=+22{,}696&F_3=+21{,}587&F_4=+19{,}140&F_5=+14{,}208&F_6=+4{,}814\\
F_{12}=-194{,}684&F_{11}=-179{,}157&F_{10}=-136{,}441&F_9=-84{,}013&F_8=-40{,}745&F_7=-12{,}117&
\end{array}
\]
\newcommand{\VfourDots}[1]{\begingroup\count255=#1\relax\loop\mathord{.}\,\advance\count255 by -1\ifnum\count255>0\repeat\endgroup}
\clearpage\SourceLeaf{507}
\par\noindent\hspace*{\parindent}Die Cosinus 4 (8 fach):
\[
\begin{array}{@{}r|r|r|r@{}}
K_0=-171{,}434&K_1=-155{,}998&K_2=-113{,}745&K_3=-62{,}426\\
K_6=+9{,}628&K_5=+2{,}091&K_4=-21{,}605&
\end{array}
\]
\par\noindent\hspace*{\parindent}Die ungeraden Cosinus 4 (16 fach):
\[M_0=-181{,}062\qquad M_1=-158{,}089\qquad M_2=-92{,}140\]
\[
\begin{array}{@{}r@{\qquad}r@{\qquad}r@{}}
M_0-2M_2=+3{,}218&M_1\VfourDots{6}2{,}19890_n&M_0+M_2=-273{,}202\\
\text{Dividirt mit 48:}&\BareRoot{768}\underline{\VfourDots{4}1{,}44268}&\text{Dividirt mit 48:}\quad-5{,}6917\\
\alpha_{12}=+0{,}067&\multicolumn{2}{r@{}}{0{,}75622_n\VfourDots{15}\text{Zahl}\ \underline{= -5{,}7045}}\\
&&\alpha_4=-11{,}396\\
&&\alpha_{20}=+0{,}013.
\end{array}
\]
\par\noindent\hspace*{\parindent}Die Cosinus 8 (16 fach):
\[P_0=-161{,}806\qquad P_1=-153{,}907\qquad P_2=-135{,}350\qquad P_3=-124{,}852\]
\[
\begin{array}{@{}r@{\qquad}r@{}}
P_0+P_3=-286{,}658&P_0-P_3=-36{,}954\\
P_1+P_2=-289{,}257&P_1-P_2=-18{,}557\\
2(P_1+P_2)=-578{,}514&P_0-P_3+P_1-P_2=-55{,}511\\
P_0+P_3+2(P_1+P_2)=-865{,}172&P_0-P_3-2(P_1-P_2)=+0{,}160\\
P_0+P_3-(P_1+P_2)=+2{,}599&\\
\text{Dividirt mit 96 und 48:}&\text{Dividirt mit 48 und 96:}\\
\alpha_0=-9{,}012&\alpha_8=-1{,}157\\
\alpha_{16}=+0{,}054&\alpha_{24}=+0{,}002.
\end{array}
\]
\par\noindent\hspace*{\parindent}Die ungeraden Cosinus 2 (8 fach):
\[H_0=+217{,}934\quad H_1=+202{,}316\quad H_2=+159{,}137\quad H_3=+105{,}600\quad H_4=+59{,}885\quad H_5=+26{,}325\]
\[
\begin{array}{@{}r@{\qquad}r@{\qquad}r@{}}
\begin{array}[t]{@{}r@{}}
H_1-H_3-H_5\VfourDots{6}1{,}84752\\
\BareRoot{1152}\VfourDots{7}\underline{\VfourDots{5}1{,}53073}\\
0{,}31679\\[4pt]
\text{Zahl}=+2{,}0739\\
\dfrac{H_0-2H_4}{48}\underline{=+2{,}0451}\\
\alpha_6=+4{,}119\\
\alpha_{18}=-0{,}029
\end{array}
&\begin{array}[t]{@{}r@{}}
H_2\VfourDots{15}2{,}20177\\
\BareRoot{768}\VfourDots{8}\underline{\VfourDots{5}1{,}44268}\\
(*)\VfourDots{16}0{,}75909\\[4pt]
H_1+H_5\VfourDots{10}2{,}35916\\
\BareRoot{1536}\VfourDots{7}\underline{\VfourDots{5}1{,}59320}\\
(**)\VfourDots{15}0{,}76596\\[4pt]
(H_1-H_5+2H_3)\VfourDots{4}2{,}58793\\
\BareRoot{4608}\VfourDots{7}\underline{\VfourDots{5}1{,}83176}\\
(***)\VfourDots{14}0{,}75617
\end{array}
&\begin{array}[t]{@{}r@{}}
\dfrac{H_0+H_4}{48}=+5{,}7879\\
(*)=+5{,}7424\\
(**)=+5{,}8339\\
(***)\underline{=+5{,}7039}\\
\alpha_2=+23{,}068\\
\alpha_{10}=+0{,}176\\
\alpha_{14}=-0{,}084\\
\alpha_{22}=-0{,}007.
\end{array}
\end{array}
\]
\par\noindent\hspace*{\parindent}Die ungeraden Cosinus selbst (4 fach):
\[D_0=+11{,}106\quad D_2=+9{,}888\quad D_4=+5{,}398\quad D_6=-6{,}020\quad D_8=-33{,}443\quad D_{10}=-66{,}957\]
\clearpage\SourceLeaf{508}
\[
\begin{array}{@{}r@{\qquad}r@{\qquad}r@{}}
\begin{array}[t]{@{}r@{}}
D_2-D_6-D_{10}\VfourDots{5}1{,}91837\\
\BareRoot{1152}\VfourDots{7}\underline{\VfourDots{5}1{,}53073}\\
0{,}38764\\[4pt]
\text{Zahl}=+2{,}4414\\
\dfrac{D_0-2D_8}{48}\underline{=+1{,}6248}\\
\tfrac12(\alpha_3+\alpha_{21})=+4{,}066\\
\tfrac12(\alpha_9+\alpha_{15})=-0{,}817
\end{array}
&\begin{array}[t]{@{}r@{}}
D_4\VfourDots{14}0{,}73223\\
\BareRoot{768}\VfourDots{8}\underline{\VfourDots{4}1{,}44268}\\
(*)\VfourDots{15}9{,}28955\\[4pt]
(D_2+D_{10})\VfourDots{8}1{,}75640_n\\
\BareRoot{1536}\VfourDots{8}\underline{\VfourDots{4}1{,}59320}\\
(**)\VfourDots{14}0{,}16320_n\\[4pt]
(D_2-D_{10}+2D_6)\VfourDots{2}1{,}81161\\
\BareRoot{4608}\VfourDots{7}\underline{\VfourDots{4}1{,}83176}\\
(***)\VfourDots{13}9{,}97985
\end{array}
&\begin{array}[t]{@{}r@{}}
\dfrac{D_0+D_8}{48}=-0{,}4654\\
(*)=+0{,}1948\\
(**)=-1{,}4561\\
(***)\underline{=+0{,}9547}\\
\tfrac12(\alpha_1+\alpha_{23})=-0{,}772\\
\tfrac12(\alpha_{11}+\alpha_{13})=+0{,}231\\
\tfrac12(\alpha_5+\alpha_{19})=-3{,}071\\
\tfrac12(\alpha_7+\alpha_{17})=+1{,}751.
\end{array}
\end{array}
\]
\par\noindent\hspace*{\parindent}Die übrigen hingegen auf folgende Art:
\[D_1=+10{,}833\quad D_3=+8{,}185\quad D_5=+0{,}986\quad D_7=-17{,}029\quad D_9=-53{,}863\quad D_{11}=-50{,}751\]
\[
\begin{array}{@{}r@{\qquad}r@{\qquad}r@{}}
\begin{array}[t]{@{}r@{}}
a\cos A\VfourDots{5}0{,}79211_n\\
\underline{a\sin A\VfourDots{5}1{,}71380}\\
A=-83^0\,10'\,15''\\[5pt]
a\VfourDots{10}1{,}71690_n\\
\underline{\BareRoot{768}\VfourDots{6}1{,}44268}\\[5pt]
22^0\,30'-A=105^0\,40'\,15''\\
\cos\VfourDots{8}9{,}43154_n\\
\sin\VfourDots{9}9{,}98355\\
\underline{\dfrac{a}{\BareRoot{768}}\VfourDots{6}0{,}27422_n}\\[5pt]
\tfrac14(\alpha_1+\alpha_7-\alpha_{17}-\alpha_{23})=+0{,}508\\
\tfrac14(\alpha_5-\alpha_{11}+\alpha_{13}-\alpha_{19})=-1{,}810
\end{array}
&\begin{array}[t]{@{}r@{}}
b\cos B\VfourDots{5}1{,}91236\\
\underline{b\sin B\VfourDots{5}1{,}76305}\\
B=35^0\,20'\,22''\\[5pt]
b\VfourDots{10}2{,}00081\\
\underline{24\VfourDots{9}1{,}38021}\\[5pt]
22^0\,30'-B=-12^0\,50'\,22''\\
\cos\VfourDots{8}9{,}98900\\
\sin\VfourDots{8}9{,}34678_n\\
\underline{\dfrac{b}{24}\VfourDots{8}0{,}62060}\\[5pt]
\tfrac12(\alpha_3-\alpha_{21})=+4{,}070\\
\tfrac12(\alpha_9-\alpha_{15})=-0{,}928
\end{array}
&\begin{array}[t]{@{}r@{}}
c\cos C\VfourDots{5}1{,}90235_n\\
\underline{c\sin C\VfourDots{5}1{,}52368_n}\\
C=22^0\,41'\,32''\\[5pt]
c\VfourDots{10}1{,}93734_n\\
\underline{48\VfourDots{9}1{,}68124}\\[5pt]
67^0\,30'-C=44^0\,48'\,28''\\
\cos\VfourDots{8}9{,}85094\\
\sin\VfourDots{8}9{,}84802\\
\underline{\dfrac{c}{48}\VfourDots{8}1{,}68124}\\[5pt]
\tfrac14(\alpha_1-\alpha_7+\alpha_{17}-\alpha_{23})=-1{,}280\\
\tfrac14(\alpha_5+\alpha_{11}-\alpha_{13}-\alpha_{19})=-1{,}271.
\end{array}
\end{array}
\]
\par\noindent Also:
\[
\begin{array}{@{}r@{\qquad}r@{}}
\tfrac12(\alpha_1-\alpha_{23})=-0{,}772&\tfrac12(\alpha_7-\alpha_{17})=+1{,}787\\
\tfrac12(\alpha_5-\alpha_{19})=-3{,}081&\tfrac12(\alpha_{11}-\alpha_{13})=+0{,}540
\end{array}
\]
\par\noindent und endlich:
\[
\begin{array}{@{}r@{\qquad}r@{\qquad}r@{}}
\alpha_1=-1{,}544&\alpha_3=+8{,}136&\alpha_5=-6{,}152\\
\alpha_7=+3{,}538&\alpha_9=-1{,}744&\alpha_{11}=+0{,}770\\
\alpha_{13}=-0{,}309&\alpha_{15}=+0{,}111&\alpha_{17}=-0{,}037\\
\alpha_{19}=+0{,}010&\alpha_{21}=-0{,}004&\alpha_{23}=-0{,}000.
\end{array}
\]
\par\noindent\hspace*{\parindent}Die Sinus (2 fach):
\[
\begin{array}{@{}r|r|r|r|r|r@{}}
C_1=+4{,}194&C_2=+8{,}072&C_3=+11{,}132&C_4=+12{,}779&C_5=+12{,}079&C_6=+7{,}499\\
C_{23}=+4{,}767&C_{22}=+9{,}378&C_{21}=+13{,}653&C_{20}=+17{,}357&C_{19}=+20{,}151&C_{18}=+21{,}511
\end{array}
\]
\[
\begin{array}{@{}r|r|r|r|r|r@{}}
C_7=-3{,}243&C_8=-22{,}962&C_9=-52{,}448&C_{10}=-83{,}309&C_{11}=-95{,}134&C_{12}=-76{,}566\\
C_{17}=+20{,}596&C_{16}=+16{,}025&C_{15}=+5{,}609&C_{14}=-13{,}518&C_{13}=-42{,}957&
\end{array}
\]
\par\noindent\hspace*{\parindent}Die Sinus 2 (4 fach):
\[
\begin{array}{@{}r|r|r|r|r|r@{}}
G_1=-0{,}573&G_2=-1{,}306&G_3=-2{,}521&G_4=-4{,}578&G_5=-8{,}072&G_6=-14{,}012\\
G_{11}=-52{,}177&G_{10}=-69{,}791&G_9=-58{,}057&G_8=-38{,}987&G_7=-23{,}839&
\end{array}
\]
\clearpage\SourceLeaf{509}
\par\noindent\hspace*{\parindent}Die Sinus 4 (8 fach):
\[
\begin{array}{@{}r|r|r@{}}
L_1=+51{,}604&L_2=+68{,}485&L_3=+55{,}536\\
L_5=+15{,}767&L_4=+34{,}409&
\end{array}
\]
\par\noindent\hspace*{\parindent}Die ungeraden Sinus 4:
\[N_1=+67{,}371\qquad N_2=+102{,}894\qquad N_3=+111{,}072\]
\[
\begin{array}{@{}r@{\qquad}r@{\qquad}r@{}}
2N_1-N_3=+23{,}670&N_2\VfourDots{8}2{,}01239&N_1+N_3=+178{,}443\\
\beta_{12}=+0{,}493&\BareRoot{768}\VfourDots{1}\underline{\VfourDots{5}1{,}44268}&\text{Dividirt mit 48:}+3{,}718\\
&\multicolumn{2}{r@{}}{0{,}56971\VfourDots{14}\text{Zahl}\ \underline{=+3{,}713}}\\
&&\beta_4=+7{,}431\\
&&\beta_{20}=+0{,}005.
\end{array}
\]
\par\noindent\hspace*{\parindent}Die Sinus 8 (16 fach):
\[
\begin{array}{@{}r@{\qquad}r@{}}
Q_1=+35{,}837&Q_1+Q_2\VfourDots{4}1{,}84456\\
Q_2=+34{,}076&Q_1-Q_2\VfourDots{4}0{,}24576\\
&\BareRoot{768}\VfourDots{1}\underline{\VfourDots{5}1{,}44268}\\
&\beta_8=+2{,}523\\
&\beta_{16}=+0{,}064.
\end{array}
\]
\par\noindent\hspace*{\parindent}Die ungeraden Sinus 2:
\[\mathrm{J}_1=-52{,}750\quad\mathrm{J}_2=-71{,}097\quad\mathrm{J}_3=-60{,}578\quad\mathrm{J}_4=-43{,}565\quad\mathrm{J}_5=-31{,}911\quad\mathrm{J}_6=-28{,}024\]
\[
\begin{array}{@{}r@{\qquad}r@{\qquad}r@{}}
\begin{array}[t]{@{}r@{}}
\mathrm{J}_1+\mathrm{J}_3-\mathrm{J}_5\VfourDots{7}1{,}91072_n\\
\BareRoot{1152}\VfourDots{7}\underline{\VfourDots{4}1{,}53073}\\
0{,}37999_n\\[4pt]
\text{Zahl}=-2{,}399\\
\dfrac{2\mathrm{J}_2-\mathrm{J}_6}{48}\underline{=-2{,}379}\\
\beta_6=-4{,}778\\
\beta_{18}=-0{,}020
\end{array}
&\begin{array}[t]{@{}r@{}}
\mathrm{J}_4\VfourDots{15}1{,}63914_n\\
\BareRoot{768}\VfourDots{8}\underline{\VfourDots{4}1{,}44268}\\
(*)\VfourDots{15}0{,}19646_n\\[4pt]
\mathrm{J}_1+\mathrm{J}_5\VfourDots{11}1{,}92768_n\\
\BareRoot{1536}\VfourDots{7}\underline{\VfourDots{4}1{,}59320}\\
(**)\VfourDots{14}0{,}33448_n\\[4pt]
\mathrm{J}_1-\mathrm{J}_5-2\mathrm{J}_3\VfourDots{6}2{,}00137\\
\BareRoot{4608}\VfourDots{7}\underline{\VfourDots{4}1{,}83176}\\
(***)\VfourDots{13}0{,}16961
\end{array}
&\begin{array}[t]{@{}r@{}}
\dfrac{\mathrm{J}_2+\mathrm{J}_6}{48}=-2{,}065\\
(*)=-1{,}572\\
(**)=-2{,}160\\
(***)\underline{=+1{,}478}\\
\beta_2=-7{,}275\\
\beta_{10}=-1{,}175\\
\beta_{14}=-0{,}189\\
\beta_{22}=-0{,}001.
\end{array}
\end{array}
\]
\par\noindent\hspace*{\parindent}Die ungeraden Sinus selbst:
\par\noindent\hspace*{\parindent}Die 2. 4. 6. 8. 10. 12 wie die ungeraden Sinus 2:
\[E_2=+17{,}450\quad E_4=+30{,}136\quad E_6=+29{,}010\quad E_8=-6{,}937\quad E_{10}=-96{,}827\quad E_{12}=-153{,}132\]
\[
\begin{array}{@{}r@{\qquad}r@{\qquad}r@{}}
\begin{array}[t]{@{}r@{}}
E_2+E_6-E_{10}\VfourDots{4}2{,}15621\\
\BareRoot{1152}\VfourDots{8}\underline{\VfourDots{4}1{,}53073}\\
0{,}62548\\[4pt]
\text{Zahl}=+4{,}222\\
\dfrac{2E_4-E_{12}}{48}\underline{=+4{,}446}\\
\tfrac12(\beta_3-\beta_{21})=+8{,}668\\
\tfrac12(\beta_9-\beta_{15})=-0{,}224
\end{array}
&\begin{array}[t]{@{}r@{}}
E_8\VfourDots{15}0{,}84117_n\\
\BareRoot{768}\VfourDots{8}\underline{\VfourDots{4}1{,}44268}\\
(*)\VfourDots{15}9{,}39849_n\\[4pt]
E_2+E_{10}\VfourDots{9}1{,}89970_n\\
\BareRoot{1536}\VfourDots{7}\underline{\VfourDots{4}1{,}59320}\\
(**)\VfourDots{14}0{,}30650_n\\[4pt]
E_2-E_{10}-2E_6\VfourDots{3}1{,}75018\\
\BareRoot{4608}\VfourDots{7}\underline{\VfourDots{4}1{,}83176}\\
(***)\VfourDots{14}9{,}91842
\end{array}
&\begin{array}[t]{@{}r@{}}
\dfrac{E_4+E_{12}}{48}=-2{,}562\\
(*)=-0{,}250\\
(**)=-2{,}025\\
(***)\underline{=+0{,}829}\\
\tfrac12(\beta_1-\beta_{23})=-5{,}666\\
\tfrac12(\beta_7-\beta_{17})=+1{,}116\\
\tfrac12(\beta_5-\beta_{19})=-3{,}508\\
\tfrac12(\beta_{11}-\beta_{13})=-0{,}042.
\end{array}
\end{array}
\]
\clearpage\SourceLeaf{510}
\par\noindent 1. 3. 5. 7. 9. 11 auf folgende Art:
\[E_1=+8{,}961\quad E_3=+24{,}785\quad E_5=+32{,}230\quad E_7=+17{,}353\quad E_9=-46{,}839\quad E_{11}=-138{,}091\]
\[
\begin{array}{@{}r@{\qquad}r@{\qquad}r@{}}
\begin{array}[t]{@{}r@{}}
a'\sin A'\VfourDots{10}0{,}92387\\
\underline{a'\cos A'\VfourDots{9}2{,}02474_n}\\
A'=-4^0\,31'\,57''\\[5pt]
a'\VfourDots{15}2{,}02610_n\\
\underline{\BareRoot{768}\VfourDots{12}1{,}44268}\\[5pt]
22^0\,30'-A'=27^0\,1'\,57''\\
\cos\VfourDots{14}9{,}94975\\
\sin\VfourDots{14}9{,}65753\\
\underline{\dfrac{a'}{\BareRoot{768}}\VfourDots{12}0{,}58342_n}\\[5pt]
\tfrac14(\beta_1-\beta_7-\beta_{17}+\beta_{23})=-3{,}413\\
\tfrac14(\beta_5+\beta_{11}+\beta_{13}+\beta_{19})=-1{,}742
\end{array}
&\begin{array}[t]{@{}r@{}}
b'\sin B'\VfourDots{10}1{,}86423\\
\underline{b'\cos B'\VfourDots{9}2{,}29027}\\
B'=20^0\,33'\,11''\\[5pt]
b'\VfourDots{15}2{,}31883\\
\underline{24\VfourDots{15}1{,}38021}\\[5pt]
B'-22^0\,30'=-1^0\,56'\,49''\\
\cos\VfourDots{14}9{,}99975\\
\sin\VfourDots{14}8{,}53116_n\\
\underline{\dfrac{b'}{24}\VfourDots{14}0{,}93862}\\[5pt]
\tfrac12(\beta_3+\beta_{21})=+8{,}677\\
\tfrac12(\beta_9+\beta_{15})=-0{,}295
\end{array}
&\begin{array}[t]{@{}r@{}}
c'\sin C'\VfourDots{10}1{,}82843_n\\
\underline{c'\cos C'\VfourDots{10}2{,}08189_n}\\
C'=29^0\,9'\,22''\\[5pt]
c'\VfourDots{15}2{,}14073_n\\
\underline{48\VfourDots{15}1{,}68124}\\[5pt]
67^0\,30'-C'=38^0\,20'\,38''\\
\cos\VfourDots{14}9{,}89449\\
\sin\VfourDots{15}9{,}79266\\
\underline{\dfrac{c'}{48}\VfourDots{2}\mathord{.}\mkern0.5mu\mathord{.}\,\VfourDots{11}0{,}45949_n}\\[5pt]
\tfrac14(\beta_1+\beta_7+\beta_{17}+\beta_{23})=-2{,}259\\
\tfrac14(\beta_5-\beta_{11}-\beta_{13}+\beta_{19})=-1{,}787.
\end{array}
\end{array}
\]
\par\noindent Also:
\[
\begin{array}{@{}r@{\qquad}r@{}}
\tfrac12(\beta_1+\beta_{23})=-5{,}672&\tfrac12(\beta_7+\beta_{17})=+1{,}154\\
\tfrac12(\beta_5+\beta_{19})=-3{,}529&\tfrac12(\beta_{11}+\beta_{13})=+0{,}045
\end{array}
\]
\par\noindent und endlich:
\[
\begin{array}{@{}r@{\qquad}r@{\qquad}r@{}}
\beta_1=-11{,}338&\beta_3=+17{,}345&\beta_5=-7{,}037\\
\beta_7=+2{,}270&\beta_9=-0{,}519&\beta_{11}=+0{,}003\\
\beta_{13}=+0{,}087&\beta_{15}=-0{,}071&\beta_{17}=+0{,}038\\
\beta_{19}=-0{,}021&\beta_{21}=+0{,}009&\beta_{23}=-0{,}006.
\end{array}
\]
\par\noindent Das ganze Resultat ist also:
\[
\begin{array}{@{}r@{\quad}r@{\qquad\qquad}r@{\quad}r@{}}
-9{,}012&&&\\
-1{,}544\cos 1&-11{,}338\sin 1&-0{,}309\cos 13&+0{,}087\sin 13\\
+23{,}068\cos 2&-7{,}275\sin 2&-0{,}084\cos 14&-0{,}189\sin 14\\
+8{,}136\cos 3&+17{,}345\sin 3&+0{,}111\cos 15&-0{,}071\sin 15\\
-11{,}396\cos 4&+7{,}431\sin 4&+0{,}054\cos 16&+0{,}064\sin 16\\
-6{,}152\cos 5&-7{,}037\sin 5&-0{,}037\cos 17&+0{,}038\sin 17\\
+4{,}119\cos 6&-4{,}778\sin 6&-0{,}029\cos 18&-0{,}020\sin 18\\
+3{,}538\cos 7&+2{,}270\sin 7&+0{,}010\cos 19&-0{,}021\sin 19\\
-1{,}157\cos 8&+2{,}523\sin 8&+0{,}013\cos 20&+0{,}005\sin 20\\
-1{,}744\cos 9&-0{,}519\sin 9&-0{,}004\cos 21&+0{,}009\sin 21\\
+0{,}176\cos 10&-1{,}175\sin 10&-0{,}007\cos 22&-0{,}001\sin 22\\
+0{,}770\cos 11&+0{,}003\sin 11&-0{,}000\cos 23&-0{,}006\sin 23\\
+0{,}067\cos 12&+0{,}493\sin 12&+0{,}002\cos 24.&
\end{array}
\]
\clearpage
\endgroup
\end{landscape}
